% =====================================================================
%  diagrams/procesos-prueba.tex
%  Figura: flujo del proceso de prueba en tres capas, con sus bucles
%          de retroalimentación.
%
%  CORRECCIONES respecto del diagrama original («Procesos de prueba.png»):
%   1. El original rotulaba «verificar correcciones» como prueba de
%      regresión. Aquí se representan por separado la prueba de
%      CONFIRMACIÓN (el defecto quedó corregido) y la de REGRESIÓN (el
%      cambio no dañó lo que ya funcionaba).
%   2. El original se leía como una secuencia lineal. Aquí el monitoreo
%      y control aparece como proceso CONCURRENTE y se dibujan los tres
%      bucles de retroalimentación: control, replanificación y mejora.
%   3. El original situaba el entorno y los datos dentro de la
%      preparación de pruebas; en la edición 2021 constituyen un proceso
%      propio (ES1/ES2), y así se representa.
%   4. El original sólo contemplaba el análisis de riesgos del proyecto;
%      la planificación abarca riesgos de producto y de proyecto.
% =====================================================================
\begin{figure}[htbp]
\centering
\begin{tikzpicture}[
  font=\scriptsize,
  proc/.style   ={draw=uvazul, thick, fill=uvazul!8, rounded corners=2pt,
                  text width=2.8cm, align=center, minimum height=1.0cm, inner sep=3pt},
  procm/.style  ={draw=polinaranja, very thick, fill=polfondo, rounded corners=2pt,
                  text width=2.9cm, align=center, minimum height=1.5cm, inner sep=3pt},
  salida/.style ={draw=gray!70, fill=gray!8, rounded corners=2pt,
                  text width=2.6cm, align=center, minimum height=0.7cm, inner sep=3pt},
  conf/.style   ={draw=instrverde, thick, fill=instrverde!8, rounded corners=2pt,
                  text width=2.7cm, align=center, minimum height=0.9cm, inner sep=3pt},
  fl/.style     ={-{Latex[length=4pt]}, draw=uvazul!85, thick},
  fb/.style     ={-{Latex[length=4pt]}, draw=polinaranja, thick, dashed},
  mej/.style    ={-{Latex[length=4pt]}, draw=instrverde, thick, dashdotted},
  et/.style     ={font=\tiny, inner sep=1pt, fill=white}
]

% ============================ CAPA ORGANIZACIONAL =====================
\node[proc] (ot) at (-4.6,0) {\textbf{Proceso organizacional}\\ OT1 desarrollar · OT2 monitorear · OT3 actualizar};
\node[salida] (pol) at (-0.6,0) {Política de prueba\\ y estrategia organizacional};
\draw[fl] (ot) -- (pol);

% ============================ CAPA DE GESTIÓN =========================
\node[proc] (tp) at (-5.6,-2.6) {\textbf{Planificación}\\ TP1--TP9\\ \emph{riesgos de producto y de proyecto}};
\node[salida] (plan) at (-5.6,-4.3) {Plan de prueba};
\draw[fl] (tp) -- (plan);

\node[procm] (tmc) at (-0.4,-3.3) {\textbf{Monitoreo y control}\\ TMC1--TMC4\\ \emph{concurrente con la ejecución}};

\node[proc] (tc) at (4.6,-3.3) {\textbf{Finalización}\\ TC1 archivar · TC2 limpiar entorno · TC3 lecciones · TC4 reportar};

\draw[fl] (pol) to[out=250,in=70] node[et,pos=0.5]{gobierna} (tp.north);

% ============================ CAPA DINÁMICA ===========================
\node[proc] (es) at (-6.0,-6.6) {\textbf{Entorno y datos}\\ ES1 establecer\\ ES2 mantener};
\node[proc] (td) at (-2.6,-6.6) {\textbf{Diseño e implementación}\\ TD1--TD6};
\node[proc] (te) at (0.8,-6.6)  {\textbf{Ejecución}\\ TE1 ejecutar · TE2 comparar · TE3 registrar};
\node[proc] (ir) at (4.2,-6.6)  {\textbf{Reporte de incidentes}\\ IR1 analizar · IR2 reportar};

\draw[fl] (es) -- (td);
\draw[fl] (td) -- (te);
\draw[fl] (te) -- node[et,pos=0.5]{no aprobado} (ir);

\draw[fl] (plan.south) to[out=270,in=120] node[et,pos=0.25]{dirige} (es.north);

% ------------- bucle interno: corrección, confirmación, regresión -----
\node[salida] (nv) at (4.2,-8.5) {Nueva versión del elemento de prueba};
\draw[fl] (ir) -- node[et,pos=0.5]{corrección} (nv);

\node[conf] (pc) at (0.3,-8.5)  {\textbf{Prueba de confirmación}\\ ¿se corrigió \emph{ese} defecto?};
\node[conf] (pr) at (0.3,-10.4) {\textbf{Prueba de regresión}\\ ¿el cambio dañó lo que ya funcionaba?};
\draw[fl] (nv.west) -- (pc.east);
\draw[fl] (nv.south) to[out=250,in=340] (pr.east);

\draw[fl] (pc.north) -- node[et,pos=0.5]{nueva ejecución} (te.south);
\draw[fl] (pr.west) to[out=180,in=250,looseness=1.2] (te.south west);

% ======================= BUCLES DE RETROALIMENTACIÓN ==================
% 1) bucle de control corto: dinámica -> monitoreo -> dinámica
\draw[fb] (te.north) to[out=100,in=280] node[et,pos=0.5]{medidas y estado} (tmc.south);
\draw[fb] (tmc.west) to[out=200,in=80] node[et,pos=0.55]{directivas de control} (td.north);

% 2) bucle de replanificación: monitoreo -> planificación
\draw[fb] (tmc.north west) to[out=150,in=20] node[et,pos=0.5]{replanificación} (tp.east);

% 3) bucle de mejora organizacional: finalización -> capa organizacional
% se rodea por encima de la capa organizacional para no atravesar ningún nodo
\draw[mej] (tc.north) |- ($(ot.north)+(0,1.25)$)
             node[et,pos=0.72]{lecciones aprendidas} -- (ot.north);

% cierre hacia finalización
\draw[fl] (ir.east) to[out=20,in=250] node[et,pos=0.5]{criterios de salida} (tc.south);

% ============================ BANDAS DE CAPA ==========================
\begin{scope}[on background layer]
  \node[fit=(ot)(pol), draw=gray!45, dashed, inner sep=7pt, rounded corners] (b1) {};
  \node[fit=(tp)(plan)(tmc)(tc), draw=gray!45, dashed, inner sep=7pt, rounded corners] (b2) {};
  \node[fit=(es)(td)(te)(ir)(nv)(pc)(pr), draw=gray!45, dashed, inner sep=7pt, rounded corners] (b3) {};
\end{scope}
\node[et, text=gray!60, anchor=west, font=\tiny\bfseries] at (b1.north west) {CAPA ORGANIZACIONAL};
\node[et, text=gray!60, anchor=west, font=\tiny\bfseries] at (b2.north west) {CAPA DE GESTIÓN DE LA PRUEBA};
\node[et, text=gray!60, anchor=west, font=\tiny\bfseries] at (b3.north west) {CAPA DINÁMICA (iterativa)};

% ============================== LEYENDA ===============================
\node[align=left, font=\scriptsize, text width=14.8cm] at (0,-12.4)
  {\textbf{Cómo leer la figura.} Las flechas azules continuas indican el flujo principal; las
   naranjas discontinuas, los bucles de control y de replanificación; la verde de trazo mixto,
   el bucle lento de mejora organizacional. El monitoreo y control \emph{no} ocupa un lugar fijo
   en la secuencia: acompaña a la ejecución de principio a fin. Los pasos de la capa dinámica se
   repiten por conjunto de pruebas, incremento o iteración hasta cumplir los criterios de salida.};

\end{tikzpicture}
\caption[Flujo del proceso de prueba con sus bucles de retroalimentación]%
{Flujo del proceso de prueba en las tres capas de \normap{2}, con el monitoreo y control
representado como proceso concurrente y con los bucles de control, replanificación y mejora.}
\label{fig:procesos}
\end{figure}

\begin{observacion}[Diagrama del flujo del proceso]
\obsproblema{El diagrama de referencia rotulaba «verificar correcciones» como prueba de
regresión, se leía como una secuencia lineal, situaba el entorno y los datos dentro de la
preparación de pruebas y contemplaba sólo el análisis de riesgos del proyecto.}
\obscambio{Confirmación y regresión aparecen como cajas distintas; el monitoreo y control se
dibuja como proceso concurrente con tres bucles de retroalimentación explícitos; entorno y datos
constituyen un proceso propio (ES1/ES2); la planificación abarca riesgos de producto y de
proyecto.}
\obsfuente{OBS §5.5 y §5.6 · PDF §2.2 · DIAG}
\end{observacion}
